% !TEX root = ../main.tex
%----------------------------------------------------------
% 06-tripartite-value.tex — Section 5: Three-Component Value
% MIGRATED from preprint A5-verification-framework.tex §E.6
%    + A6-plain-language-verification.tex three-lanes section
% Removed: Clyfar-specific language, appendix framing
% Ozone categories → SPC categories
%----------------------------------------------------------

\section{Three-Component Value}\label{sec:tripartite_value}

Each component of the five-number scorecard
(Table~\ref{tab:three_to_five}) should earn its place by providing
diagnostic information the others cannot.  This section first examines
the independent value of the three components---possibility (shape
quality), ignorance (system confidence), and conditional necessity
(dominance)---and then steps back to show how the native possibilistic
facet relates to probabilistic and categorical verification.

%----------------------------------------------------------
\subsection{Independent Value of Each Component}\label{sec:component_value}

\paragraph{Possibility: does the shape point at the truth?}

The possibility component ($\alpha^*$, $\eta$, $\delta$) captures
whether the system's distributional shape favours the observed outcome.
Its value reduces to a familiar verification question: does the system's
support margin $\delta$ exceed that of a na\"ive baseline?  Comparing
$\delta$ against climatology or persistence would be the possibilistic
analogue of comparing Brier skill scores
(Section~\ref{sec:limitations} notes that formalising this comparison
requires a possibilistic climatology that remains an open problem).  For SPC outlooks, climatology
peaks at the most frequent category and assigns near-zero possibility to
rare outcomes like \spc{HIGH}; a skilful possibilistic system should
outperform this baseline on days when \spc{ENH} or higher categories
verify, exhibiting positive $\delta$ across the verification sample.

\paragraph{Ignorance: does self-reported uncertainty track difficulty?}

The ignorance component ($\ign$) carries forward from the pre-event
framework as a forecast property rather than a verification metric
(Table~\ref{tab:three_to_five}), so its diagnostic value requires
explicit justification.  If high-\ign{} forecasts tend to have lower
mean $\delta$ than low-\ign{} forecasts, then $\ign$ tracks actual
forecast difficulty---the possibilistic analogue of the
\textit{spread--skill relationship} in ensemble forecasting, where high
ensemble spread should correspond to higher forecast error.

The implications are decisive for system design.  If \ign{} does
\emph{not} track difficulty, the subnormality signal is uninformative
and normalisation is costless---one might as well enforce
$\max(\pi) = 1$ and simplify the framework.  If \ign{} \emph{does}
track difficulty, normalisation destroys operationally valuable
information: a forecaster who normalises away high-\ign{} signals loses
the ability to distinguish confident forecasts from hedged ones.

\paragraph{Conditional necessity: does claimed certainty verify?}

When the system assigns high conditional necessity to an outcome, that
outcome should verify more often than the base rate.  This is the
possibilistic analogue of \textit{reliability conditional on
confidence}: ensemble forecasters routinely ask ``when the ensemble
gives ${>}80\%$ chance of severe weather, does it verify ${>}80\%$ of
the time?''  Here the question becomes: when $\Nc(\hat{c}) > \tau$ for
threshold $\tau$, does the conditional hit rate exceed climatology?
(Note the shift: here we condition on the forecast's own peak category
$\hat{c}$ rather than the observation, since reliability is a
pre-verification property of the issued forecast.)
If so, high-necessity forecasts identify genuinely confident
predictions---not just lucky ones.

Figure~\ref{fig:reliability} tests this by plotting the conditional hit
rate against the $\Nc(\hat{c})$ threshold~$\tau$: at each $\tau$, only
forecasts whose peak category has $\Nc(\hat{c}) \geq \tau$ are included.
The curve is nearly flat because the synthetic system already has a high
baseline accuracy (${\sim}83\%$).  The modest tick-up at high
thresholds ($\tau > 0.8$) is notable: it confirms that the most
confident forecasts---where the peak category strongly dominates all
alternatives---verify even more reliably, so $\Nc(\hat{c})$ carries a genuine
confidence signal.  With real data or a weaker baseline system, this
separation would be more pronounced.

\begin{figure}[t]
\centering
\includegraphics[width=0.75\columnwidth]{reliability_curves.png}
\caption{Conditional hit rate vs.\ $\Nc(\hat{c})$ threshold for the
    synthetic reforecast.  At each $\tau$, only forecasts with
    $\Nc(\hat{c}) \geq \tau$ are retained; the purple curve shows
    their hit rate.  Baselines: random chance ($1/K$, dashed) and
    unconditional system accuracy (dot-dashed).  Shaded band: notional
    confidence interval.  Inset: declining sample size at high $\tau$.
    The near-flat curve reflects the synthetic system's high baseline
    accuracy; the uptick at $\tau > 0.8$ confirms that $\Nc(\hat{c})$ adds a
    genuine confidence signal.}
\label{fig:reliability}
\end{figure}

%----------------------------------------------------------
\subsection{Three Verification Facets}\label{sec:three_lanes}

The paper proposes three expressions for forecast uncertainty and five
metrics to evaluate forecast quality.  No single facet captures the full
picture, because each answers a different question and each has a blind
spot the others fill.

The \textbf{categorical facet} reduces each distribution to its peak
category $\hat{c} = \arg\max \pi(\omega)$ (ties broken by severity)
and applies threshold-based contingency scores.  For each severity
threshold~$t$ (\spc{MRGL}+, \spc{SLGT}+, \ldots, \spc{HIGH}), the
forecast issues ``yes'' if $\hat{c} \geq t$ and ``no'' otherwise; the
observation is ``yes'' if $c_{\mathrm{obs}} \geq t$.  Four scores
evaluate each threshold:
\begin{align*}
\mathrm{POD}  &= \frac{\text{hits}}{\text{hits} + \text{misses}}
    &&\text{(did the system detect events?)} \\
\mathrm{FAR}  &= \frac{\text{false alarms}}{\text{hits} + \text{false alarms}}
    &&\text{(when it said ``yes,'' was it right?)} \\
\mathrm{CSI}  &= \frac{\text{hits}}{\text{hits} + \text{misses} + \text{false alarms}}
    &&\text{(joint penalty for misses and false alarms)}
\end{align*}
POD alone rewards over-forecasting; FAR alone rewards
under-forecasting; CSI balances both.  The Heidke skill score (HSS),
computed from the full $K \times K$ confusion matrix, measures
agreement beyond what random chance would produce.
Because verification is threshold-based, a peak forecast of \spc{MDT}
when \spc{HIGH} verifies registers as a hit at all thresholds up to
\spc{MDT}+ and a miss only at \spc{HIGH}---partial credit for
near-misses in severity.
The categorical facet asks ``did the peak-category forecast identify the
correct severity category?'' and discards all distributional
information: a correct peak says nothing about sharpness, hedging, or
confidence.

The \textbf{probabilistic facet} (IG via probability conversion) asks
``was the probabilistic forecast informative?''  The conversion is
lossy---it compresses the full possibility shape into $n{+}1$
probabilities, with subnormality collapsed into a single ignorance bin.

The \textbf{native possibilistic facet} (five-number scorecard) asks
``was the possibility distribution well-shaped, sharp, and honest?''
It has no connection to probabilistic skill: a well-shaped distribution
might still give poor probabilities after conversion.

Together, the three facets can expose failures that any single facet may miss.
A forecast with correct peak category but flat possibility and poor
IG is a broken clock that happens to be right---the categorical facet
passes but the other two fail.  A forecast with strong IG but high
ignorance and low $\Nc^*$ got lucky after conversion---the raw system
was hedging.  A forecast with sharp possibility and high $\Nc^*$ but
poor IG may indicate distortion introduced by the conversion---the
probability derivation is compressing or shifting the signal.

In practice, a verification analyst would compile scorecard aggregates
across a test sample, tracking changes between model versions or
seasons.  Figure~\ref{fig:scorecard_table} illustrates this convention
using synthetic data for four model versions.  Metrics are grouped by
verification facet; filled triangles denote statistically significant
change (hypothetical; significance assessed by paired permutation test
at the 5\% level) and marker size encodes magnitude, following the
scorecard style used by ECMWF and WMO Lead Centres for NWP
inter-comparisons.

\textit{Reading the scorecard at a glance:} a column of large green
triangles means the new version is broadly better; a red triangle amid
green ones flags a trade-off that warrants investigation.  The example
tells a common development story: version~2.0 improved discrimination
but increased ignorance (the system learned to hedge more honestly);
version~3.0 delivered broad gains across all three facets; version~4.0
pushed discrimination further at the cost of slight reliability
degradation---a trade-off the three-facet view makes immediately visible
but that a single-metric evaluation would likely miss.

\begin{figure}[t]
\centering
\includegraphics[width=\columnwidth]{scorecard_table.png}
\caption{Verification scorecard for a hypothetical possibilistic
    forecasting system across three development versions, each compared
    to a v1.0 baseline.  Rows are grouped into three verification facets:
    native possibilistic (five-number scorecard), probabilistic
    (information gain and its decomposition),
    and categorical (POD, FAR, CSI, HSS).  Filled triangles indicate
    statistically significant improvement (green, upward) or degradation
    (red, downward); open triangles indicate non-significant changes.
    Marker size encodes the magnitude of the change.  Note the
    $\overline{H_\Pi}$ degradation in v2.0: the system became more
    ignorant before improving in subsequent versions---a signature of
    recalibrated uncertainty that the native possibilistic facet detects
    but the probabilistic facet may absorb.
    Ignorance changes (amber) are context-dependent: increased $\ign$
    may indicate honest self-awareness or declining skill, depending on
    whether discrimination also changes.}
\label{fig:scorecard_table}
\end{figure}

\FloatBarrier
